\documentclass[10pt,a4paper]{amsart}
\usepackage{amsmath,amssymb,amsthm,graphicx}
% Retained preamble for the Stacks Project chapter snapshot in this directory.
% The chapter's own preamble.tex is a copy of the shared Stacks Project preamble; the block
% below is sliced MECHANICALLY, line for line, from it. Nothing was added, removed or
% redefined; the wrapper class, the xr-hyper/hyperref wiring and the externaldocument list
% are omitted because the assembler supplies the document class.
%
% lines taken from preamble.tex: 3, 5, 9, 11, 12, 156, 157, 158, ... (51 lines total)
%
% package coverage restored from the chapter's own preamble (2026-10-13): xy, xy-2cell

\usepackage[all]{xy}
\xyoption{2cell}
\UseAllTwocells

\newenvironment{reference}{\comment}{\endcomment}
\newenvironment{slogan}{\comment}{\endcomment}
\newenvironment{history}{\comment}{\endcomment}
\theoremstyle{plain}
\newtheorem{theorem}[subsection]{Theorem}
\newtheorem{proposition}[subsection]{Proposition}
\newtheorem{lemma}[subsection]{Lemma}
\theoremstyle{definition}
\newtheorem{definition}[subsection]{Definition}
\newtheorem{example}[subsection]{Example}
\newtheorem{exercise}[subsection]{Exercise}
\newtheorem{situation}[subsection]{Situation}
\theoremstyle{remark}
\newtheorem{remark}[subsection]{Remark}
\newtheorem{remarks}[subsection]{Remarks}
\numberwithin{equation}{subsection}
\def\lim{\mathop{\mathrm{lim}}\nolimits}
\def\colim{\mathop{\mathrm{colim}}\nolimits}
\def\Spec{\mathop{\mathrm{Spec}}}
\def\Hom{\mathop{\mathrm{Hom}}\nolimits}
\def\Ext{\mathop{\mathrm{Ext}}\nolimits}
\def\SheafHom{\mathop{\mathcal{H}\!\mathit{om}}\nolimits}
\def\SheafExt{\mathop{\mathcal{E}\!\mathit{xt}}\nolimits}
\def\Sch{\mathit{Sch}}
\def\Mor{\mathop{\mathrm{Mor}}\nolimits}
\def\Ob{\mathop{\mathrm{Ob}}\nolimits}
\def\Sh{\mathop{\mathit{Sh}}\nolimits}
\def\NL{\mathop{N\!L}\nolimits}
\def\CH{\mathop{\mathrm{CH}}\nolimits}
\def\proetale{{pro\text{-}\acute{e}tale}}
\def\etale{{\acute{e}tale}}
\def\QCoh{\mathit{QCoh}}
\def\Ker{\mathop{\mathrm{Ker}}}
\def\Im{\mathop{\mathrm{Im}}}
\def\Coker{\mathop{\mathrm{Coker}}}
\def\Coim{\mathop{\mathrm{Coim}}}
\DeclareMathSymbol{\boxtimes}{\mathbin}{AMSa}{"02}
\def\QCohstack{\mathcal{QC}\!\mathit{oh}}
\def\Cohstack{\mathcal{C}\!\mathit{oh}}
\def\Spacesstack{\mathcal{S}\!\mathit{paces}}
\def\Quotfunctor{\mathrm{Quot}}
\def\Hilbfunctor{\mathrm{Hilb}}
\def\Curvesstack{\mathcal{C}\!\mathit{urves}}
\def\Polarizedstack{\mathcal{P}\!\mathit{olarized}}
\def\Complexesstack{\mathcal{C}\!\mathit{omplexes}}
\def\Pic{\mathop{\mathrm{Pic}}\nolimits}
\def\Picardstack{\mathcal{P}\!\mathit{ic}}
\def\Picardfunctor{\mathrm{Pic}}
\def\Deformationcategory{\mathcal{D}\!\mathit{ef}}
\usepackage{cleveref}
\crefname{originallemma}{original lemma}{original lemmas}
\Crefname{originallemma}{Original lemma}{Original lemmas}
\makeatletter
\newlabel{morphisms-lemma-chevalley}{{054J}{source}{}{sourcepointer-1}{}}
\newlabel{morphisms-lemma-chevalley@cref}{{[originallemma][0][]054J}{[1][0][]source}{}{}{}}
\makeatother
\crefname{originallemma}{original lemma}{original lemmas}
\Crefname{originallemma}{Original lemma}{Original lemmas}
\makeatletter
\newlabel{smoothing-section-singular-ideal}{{07C4}{source}{}{sourcepointer-2}{}}
\newlabel{smoothing-section-singular-ideal@cref}{{[originallemma][0][]07C4}{[1][0][]source}{}{}{}}
\makeatother
\crefname{originallemma}{original lemma}{original lemmas}
\Crefname{originallemma}{Original lemma}{Original lemmas}
\makeatletter
\newlabel{topology-lemma-generic-point-in-constructible}{{005K}{source}{}{sourcepointer-3}{}}
\newlabel{topology-lemma-generic-point-in-constructible@cref}{{[originallemma][0][]005K}{[1][0][]source}{}{}{}}
\makeatother
\title{Stacks, Examples of Deformations: lemma-jouanolou-type-thing}
\author{The Stacks Project Authors}
\date{stacks-project:examples-defos.tex:local-snapshot-20261012; GNU FDL 1.2 or later}
\begin{document}
\maketitle
\noindent\textit{Editorial note.} Statement span is lines 3150-3166 of sources/stacks-examples-defos/pinned-original-source.tex; proof span is lines 3168-3213, the complete author proof (44 body lines).\par\medskip
\noindent\textit{Editorial note.} Navigation stubs: 3 cited-but-not-included labels are supplied as non-load-bearing pointers with their pinned tags where the snapshot tag index has them.\par\medskip
\noindent\textit{Editorial note.} extra\_packages is empty on purpose: the retained preamble already loads xy with options and a second declaration clashes in the assembler wrapper.\par\medskip
\noindent\textit{Editorial note.} Text is retained verbatim, including the author's own typos and placeholders.\par\medskip
\noindent\textit{Editorial note.} Queue index 9 of queue\_examples-defos.json (pb=1896, stmt 3150-3166, proof 3168-3213).\par\medskip
\noindent\textit{Allowed standard prerequisites (editorial).}\begin{itemize}
\item Explicit prerequisite boundary. The cited inputs are prerequisites; the argument in the retained span is not replaced by any cited result.
\end{itemize}\medskip
\ifdefined\hypertarget\hypertarget{sourcepointer-1}{}\fi
\noindent\textit{Source reference.} Original lemma 054J: sources/stacks-examples-defos/pinned-original-source.tex, source lines 3213--3213. Cited by the retained statement or proof but not part of the collected spans; supplied only as a navigation pointer. The argument does not depend on its content.\par\medskip
\ifdefined\hypertarget\hypertarget{sourcepointer-2}{}\fi
\noindent\textit{Source reference.} Original lemma 07C4: sources/stacks-examples-defos/pinned-original-source.tex, source lines 3213--3213. Cited by the retained statement or proof but not part of the collected spans; supplied only as a navigation pointer. The argument does not depend on its content.\par\medskip
\ifdefined\hypertarget\hypertarget{sourcepointer-3}{}\fi
\noindent\textit{Source reference.} Original lemma 005K: sources/stacks-examples-defos/pinned-original-source.tex, source lines 3213--3213. Cited by the retained statement or proof but not part of the collected spans; supplied only as a navigation pointer. The argument does not depend on its content.\par\medskip
\begin{lemma}
\label{lemma-jouanolou-type-thing}
Let $k$ be a field. Let $1 \leq c \leq n$ be integers.
Let $f_1, \ldots, f_c \in k[x_1, \ldots x_n]$ be elements.
Let $a_{ij}$, $0 \leq i \leq n$, $1 \leq j \leq c$ be
variables. Consider
$$
g_j = f_j + a_{0j} + a_{1j}x_1 + \ldots + a_{nj}x_n \in
k[a_{ij}][x_1, \ldots, x_n]
$$
Denote $Y \subset \mathbf{A}^{n + c(n + 1)}_k$
the closed subscheme cut out by $g_1, \ldots, g_c$.
Denote $\pi : Y \to \mathbf{A}^{c(n + 1)}_k$ the projection
onto the affine space with variables $a_{ij}$.
Then there is a nonempty Zariski open 
of $\mathbf{A}^{c(n + 1)}_k$ over which $\pi$ is smooth.
\end{lemma}

\begin{proof}
Recall that the set of points where $\pi$ is smooth is open.
Thus the complement, i.e., the singular locus, is closed.
By Chevalley's theorem (in the form of
Morphisms, Lemma \ref{morphisms-lemma-chevalley})
the image of the singular locus is constructible.
Hence if the generic point of $\mathbf{A}^{c(n + 1)}_k$
is not in the image of the singular locus, then
the lemma follows (by Topology, Lemma
\ref{topology-lemma-generic-point-in-constructible} for example).
Thus we have to show there is no point
$y \in Y$ where $\pi$ is not smooth mapping to
the generic point of $\mathbf{A}^{c(n + 1)}_k$.
Consider the matrix of partial derivatives
$$
(\frac{\partial g_j}{\partial x_i}) =
(\frac{\partial f_j}{\partial x_i} + a_{ij})
$$
The image of this matrix in $\kappa(y)$ must have rank $< c$
since otherwise $\pi$ would be smooth at $y$, see discussion in
Smoothing Ring Maps, Section \ref{smoothing-section-singular-ideal}.
Thus we can find $\lambda_1, \ldots, \lambda_c \in \kappa(y)$
not all zero such that the vector $(\lambda_1, \ldots, \lambda_c)$
is in the kernel of this matrix.
After renumbering we may assume $\lambda_1 \not = 0$.
Dividing by $\lambda_1$ we may assume our vector has
the form $(1, \lambda_2, \ldots, \lambda_c)$.
Then we obtain
$$
a_{i1} = -
\frac{\partial f_j}{\partial x_1} -
\sum\nolimits_{j = 2, \ldots, c} \lambda_j(\frac{\partial f_j}{\partial x_i} + a_{ij})
$$
in $\kappa(y)$ for $i = 1, \ldots, n$. Moreover, since $y \in Y$ we also
have
$$
a_{0j} = -f_j - a_{1j}x_1 - \ldots - a_{nj}x_n
$$
in $\kappa(y)$. This means that the subfield of $\kappa(y)$
generated by $a_{ij}$ is contained in the subfield of $\kappa(y)$
generated by the images of $x_1, \ldots, x_n, \lambda_2, \ldots, \lambda_c$,
and $a_{ij}$ except for $a_{i1}$ and $a_{0j}$.
We count and we see that the transcendence degree of this is
at most $c(n + 1) - 1$. Hence $y$ cannot map to the generic point
as desired.
\end{proof}

\end{document}
